% sec_introduccion.tex
% Section 1: Introduction
\section{Introduction}

\IEEEPARstart{I}{n} the microstructure of electronic financial markets, processing latency constitutes the key differentiating factor between capturing or losing statistical arbitrage opportunities. In High-Frequency Trading (HFT) systems, where opportunity windows are measured in nanoseconds, any source of non-determinism in the execution critical path represents a quantifiable operational risk \cite{bilokon2025}.

Modern processors based on the x86-64 architecture implement 14- to 19-stage super-scalar pipelines with speculative execution. Upon encountering each conditional jump instruction (\texttt{jcc}), the Branch Prediction Unit (BPU) --- typically a TAGE (\textit{Tagged Geometric History Length}) predictor using 2-bit saturating counters --- speculates on the branch outcome before the condition is formally evaluated \cite{intel2012}. A branch misprediction triggers a complete pipeline flush, imposing a deterministic penalty of 13--18 clock cycles \cite{fog2023}.

Bilokon et al. \cite{bilokon2025} propose an architectural solution: \textbf{semi-static conditions}. This technique completely eliminates the conditional jump instruction from the \textit{hot path} using Self-Modifying Code (SMC) or indirect branches, delegating state updates to an asynchronous \textit{cold path}.

The present work pursues four main objectives:
\begin{enumerate}
    \item Quantify the real \textit{speedup} of the semi-static technique over standard \texttt{switch}/\texttt{if-else} control blocks under uniform-random conditions.
    \item Measure the runtime \textit{overhead} of \texttt{std::function} compared to raw function pointers.
    \item Evaluate the multi-threading viability of this pattern.
    \item Identify the boundary conditions where this technique becomes counterproductive.
\end{enumerate}
